% Deutsche Parallelfassung zu alignment_is_not_a_curve.tex

\section{Alignment}

\subsection{Motivation}

In der klassischen Geometrie wird ein Eisenbahngleis häufig als Kurve
\(\gamma(s)\) im Raum repräsentiert.

Diese Repräsentation reicht für Ingenieurzwecke nicht aus, weil sie die
intrinsischen Größen nicht erfasst, die Konstruktion, Dynamik und Realisierung
bestimmen.

Insbesondere sind Krümmung und Überhöhung keine abgeleiteten Eigenschaften,
sondern primäre Entwurfsvariablen.

\subsection{Definition}

\begin{definition}[Alignment]
Ein \textbf{Alignment} ist ein über der Bogenlänge
\(s \in [0,L]\) definiertes strukturiertes Objekt mit folgenden Komponenten:

\begin{itemize}
\item einer Krümmungsfunktion
\[
\kappa : [0,L] \rightarrow \mathbb{R},
\]

\item einer Überhöhungsfunktion
\[
\phi : [0,L] \rightarrow \mathbb{R},
\]

\item einer Anfangspose
\[
(\gamma_0, \theta_0),
\]
\end{itemize}

zusammen mit den induzierten Zustandsvariablen
\[
\theta(s) = \theta_0 + \int_0^s \kappa(\sigma)\,\dd\sigma,
\]
\[
\gamma(s) = \gamma_0 + \int_0^s
(\cos\theta(\sigma), \sin\theta(\sigma))\,\dd\sigma.
\]
\end{definition}

\subsection{Interpretation}

Ein Alignment wird nicht durch seine Position \(\gamma(s)\), sondern durch
seine Krümmung \(\kappa(s)\) definiert.

Die Kurve \(\gamma(s)\) ist eine durch Integration gewonnene abgeleitete
Größe.

Krümmung ist somit die primäre Entwurfsvariable, während Position eine Folge
ist.

\subsection{Zustandsstruktur}

\begin{definition}[Alignment-Zustand]
Der \textbf{Zustand} eines Alignments an der Position \(s\) ist gegeben durch
\[
(\gamma(s), \theta(s), \kappa(s), \phi(s)).
\]
\end{definition}

Dieser Zustand koppelt Geometrie, Orientierung und Dynamik in einer
einheitlichen Repräsentation.

\subsection{Strukturelle Natur}

\begin{proposition}
Ein Alignment ist ein niedrigdimensionales strukturiertes Objekt, auch wenn
seine Realisierung hochdimensionale Daten umfassen kann.
\end{proposition}

Diese Struktur ermöglicht:
\begin{itemize}
\item effiziente Speicherung,
\item robuste Optimierung,
\item sinnvolle Interpretation.
\end{itemize}

\subsection{Beziehung zur Realisierung}

Das Alignment repräsentiert die beabsichtigte Geometrie.

Das physisch gebaute Gleis ist durch den Realisierungsoperator gegeben:
\[
\gamma_{\mathrm{real}} = \mathcal{R}(\kappa, \phi).
\]

Das Alignment existiert damit im Entwurfsraum, während das Gleis im
physischen Raum existiert.

\subsection{Beziehung zu Daten}

Messdaten liefern üblicherweise abgetastete Positionen oder Polylinien.

Ein Alignment muss aus solchen Daten über inverse Probleme rekonstruiert
werden.

\subsection{Zentrale Erkenntnis}

\begin{proposition}
Ein Alignment ist keine Kurve, sondern ein krümmungsgetriebenes generatives
Modell einer Kurve.
\end{proposition}

\subsection{Zusammenfassung}

\begin{summary}
Das Alignment ist das zentrale Objekt der Eisenbahngeometrie.

Es wird durch Krümmung und Überhöhung definiert, aus denen alle geometrischen
und dynamischen Größen abgeleitet werden.

Diese Sicht verschiebt den Schwerpunkt von der Beschreibung von Kurven zur
Definition generativer Strukturen und bildet die Grundlage des
Alignment-Based Information Modelling.
\end{summary}
